% !TEX root = ../main.tex

\chapter{Conclusioni}
\label{chap:conclusions}

Questa tesi ha affrontato il problema del rilevamento e dell’analisi di \textit{New Psychoactive Substances} in contesti operativi complessi e dinamici, proponendo una soluzione orientata al supporto investigativo sul campo attraverso la progettazione e la validazione di un laboratorio mobile simulato.
Il lavoro si inserisce nel contesto del progetto \textsc{PENTION} e si concentra in particolare sulla definizione, realizzazione e valutazione dell’artefatto \textsc{PENTION-M}, concepito come sistema cyber–fisico basato su simulazione fisica, \textit{Physics-Informed Machine Learning} e pipeline \textit{MLOps} orientate alla tracciabilità forense.

L’obiettivo principale della tesi non è stato quello di sviluppare un singolo modello predittivo, ma di progettare un’architettura coerente e verificabile capace di integrare modelli fisici, algoritmi di apprendimento automatico e requisiti operativi reali all’interno di un unico flusso end-to-end.
In questo senso, il contributo del lavoro va oltre la dimensione puramente tecnica, collocandosi all’intersezione tra ricerca applicata, ingegneria del software e supporto alle attività investigative.

\section{Sintesi dei contributi}

Il primo contributo della tesi riguarda la progettazione di un’architettura cyber–fisica modulare per il rilevamento mobile di sostanze chimiche, strutturata in otto layer funzionali.
Tale architettura consente di orchestrare la simulazione dei sensori, l’ingestion e la validazione dei dati, l’inferenza dei modelli fisici e PIML, il monitoring operativo e la generazione di evidenze forensi verificabili.
La scelta di una struttura a microservizi containerizzati ha permesso di garantire modularità, estendibilità e isolamento delle componenti, aspetti fondamentali in un contesto operativo reale.

Un secondo contributo rilevante è rappresentato dall’integrazione di modelli di \textit{Physics-Informed Machine Learning} all’interno della pipeline.
In particolare, il lavoro ha mostrato come l’introduzione di vincoli fisici espliciti nei modelli di correzione della dispersione e di localizzazione della sorgente consenta di migliorare la coerenza fisica degli output, riducendo la produzione di risultati non plausibili dal punto di vista del dominio.
Questo approccio dimostra il valore del PIML come strumento di compromesso tra fedeltà fisica e capacità predittiva in contesti caratterizzati da dati incompleti o rumorosi.

Un ulteriore contributo è rappresentato dallo sviluppo e dalla validazione di un modulo di classificazione delle NPS basato su un dataset EI esteso e integrato da più fonti.
La costruzione del dataset \textit{PENTION\_EI\_Complete} e il successivo retraining dei modelli hanno consentito di ottenere un classificatore più robusto e generalizzabile rispetto alle versioni precedenti, introducendo al contempo meccanismi di calibrazione probabilistica adattiva per il controllo dell’overconfidence, in linea con i requisiti operativi delle forze dell’ordine.

Infine, la progettazione di una pipeline \textit{MLOps} orientata alla tracciabilità forense costituisce un contributo metodologico significativo.
L’introduzione di logging strutturato, versioning dei modelli, monitoring delle metriche runtime e generazione di bundle forensi firmati consente di garantire riproducibilità, auditabilità e trasparenza del processo decisionale, aspetti fondamentali per l’utilizzo di sistemi basati su intelligenza artificiale in ambito investigativo.

\section{Impatto scientifico}

Dal punto di vista scientifico, la tesi contribuisce alla letteratura sui sistemi cyber–fisici e sul \textit{Physics-Informed Machine Learning} mostrando un’applicazione concreta di tali paradigmi in un dominio forense.
Il lavoro evidenzia come la combinazione di modelli fisici e apprendimento automatico possa essere utilizzata non solo per migliorare l’accuratezza predittiva, ma anche per aumentare la spiegabilità e la coerenza dei risultati, elementi spesso trascurati nei modelli puramente data-driven.

Inoltre, l’adozione del paradigma \textit{Design Science Research} ha permesso di strutturare il lavoro come un processo di ricerca iterativo, guidato dai requisiti degli stakeholder e validato attraverso feedback progressivi.
Questo approccio rafforza il valore metodologico della tesi, dimostrando come la DSR possa essere efficacemente applicata allo sviluppo di sistemi complessi in ambiti ad alta criticità come quello forense.

\section{Impatto operativo}

Dal punto di vista operativo, \textsc{PENTION-M} rappresenta un prototipo realistico di laboratorio mobile capace di supportare le attività di perlustrazione e rilevamento precoce delle forze dell’ordine.
Pur operando in un contesto simulato, il sistema riflette in modo fedele i workflow investigativi reali, integrando misurazioni ambientali, modelli di dispersione, analisi predittive e strumenti di visualizzazione in tempo reale.

L’attenzione alla tracciabilità, alla riproducibilità e alla spiegabilità dei risultati risponde direttamente alle esigenze espresse dagli stakeholder coinvolti nelle attività di Requirements Engineering.
In questo senso, il lavoro fornisce una base solida per futuri sviluppi orientati all’adozione operativa, riducendo il gap tra prototipazione accademica e utilizzo reale sul campo.

\section{Considerazioni finali}

Nel suo complesso, questa tesi dimostra come sia possibile progettare e validare un sistema complesso per il supporto investigativo senza ricorrere esclusivamente a dati reali o a sperimentazioni sul campo, sfruttando invece un approccio simulativo rigoroso e metodologicamente fondato.
La validazione \textit{by construction}, integrata nel ciclo di sviluppo secondo il paradigma DSR, ha consentito di valutare l’artefatto in modo coerente con i suoi obiettivi e con i vincoli del contesto applicativo.

\textsc{PENTION-M} non va interpretato come un prodotto finito, ma come una piattaforma di ricerca e sperimentazione che pone le basi per futuri studi sul campo, integrazioni hardware e validazioni operative.
Il valore principale del lavoro risiede nella dimostrazione che sistemi di intelligenza artificiale per il dominio forense possono essere progettati in modo trasparente, tracciabile e scientificamente rigoroso, offrendo un supporto concreto alle attività investigative senza sostituirsi al giudizio umano.

In questo senso, la tesi contribuisce sia allo sviluppo del progetto \textsc{PENTION} sia alla più ampia riflessione su come integrare intelligenza artificiale, fisica e ingegneria del software in contesti critici per la sicurezza pubblica.
